%============================================================
% PAKET-PAKET LATEX (khusus engine XeLaTeX)
% Compile dengan: latexmk -xelatex main.tex
% Font Times New Roman diambil dari font sistem
% (tersedia di Windows/macOS dan Overleaf).
%============================================================

% ── Font & dasar ─────────────────────────────────────────────
\usepackage{fontspec}                 % pengganti inputenc/fontenc di XeLaTeX
\setmainfont{Times New Roman}         % sesuai ketentuan kampus
\usepackage[a4paper,
            left=4cm, right=3cm,      % penjilidan: kiri 4 cm
            top=4cm,  bottom=3cm]{geometry}
\usepackage{setspace}                 % \singlespacing, \onehalfspacing, \doublespacing
\usepackage{indentfirst}              % paragraf pertama juga menjorok
\usepackage{graphicx}
\usepackage{xcolor}                   % dipakai untuk teks tersembunyi (putih)
\usepackage{array}
\usepackage{longtable}                % tabel multi-halaman dengan "(Lanjutan)"
\usepackage{enumitem}
\usepackage{amsmath}                  % rumus matematika + \eqref

% ── Judul bab/sub-bab & daftar isi ───────────────────────────
\usepackage{titlesec}
\usepackage[titles]{tocloft}          % opsi `titles`: biar tidak bentrok titlesec

% ── Caption tabel & gambar ───────────────────────────────────
\usepackage{caption}

% ── Header/footer & halaman ──────────────────────────────────
\usepackage{fancyhdr}
\usepackage{pdfpages}                 % memasukkan scan (kartu bimbingan) di lampiran
\usepackage{etoolbox}

% ── Bahasa & bibliografi ─────────────────────────────────────
\usepackage[indonesian]{babel}
\usepackage{csquotes}
\usepackage[%
  backend=biber,
  style=authoryear,
  sorting=nyt,          % urut nama (belakang dulu), tahun, judul — alfabetis
  giveninits=true,      % nama depan ditulis inisial: Wibowo, A. R.
  maxbibnames=2,        % 1–2 pengarang di daftar pustaka: tulis semua;
  minbibnames=1,        % 3+ pengarang : pengarang pertama + "dkk."
  maxnames=2,           % batas yang sama untuk sitasi dalam teks
  minnames=1,
  uniquename=false,
  uniquelist=false,
  dashed=false          % pengarang berulang tetap ditulis penuh
]{biblatex}
\addbibresource{refs/references.bib}

% ── Hyperlink (selalu terakhir) ──────────────────────────────
\usepackage[hidelinks]{hyperref}
